\documentclass[10pt,a4paper]{amsart}
\usepackage[margin=19mm]{geometry}
\usepackage{amsmath,amssymb,mathtools}
\usepackage[scr=rsfs,cal=euler]{mathalfa}
\usepackage{xfrac}
\usepackage[unicode,hidelinks,bookmarks=false]{hyperref}
\newcommand{\ie}{i.e.\ }
\newcommand{\cf}{cf.\ }
\newcommand{\resp}{respectively\ }
\newcommand{\Subsection}{\subsection}
\providecommand{\nobreakdash}{\nobreak\hbox{-}}
\setlength{\emergencystretch}{2em}
\multlinegap0pt
\usepackage{amssymb}
\usepackage[shortlabels]{enumitem}
\usepackage{float}
\usepackage{mathtools}
\newtheorem{proposition}{Proposition}
\title{Second-Jet Equivariant $\eta$ Separations on Lens Spaces}
\author{Sanchita Sharma}
\date{Source: arXiv:2606.06403v1 (CC BY 4.0)}
\begin{document}
\maketitle

\noindent\emph{Source and licence.} Second-Jet Equivariant $\eta$ Separations on Lens Spaces. Version arXiv:2606.06403v1 (CC BY 4.0), distributed by arXiv under the Creative Commons Attribution 4.0 International licence, http://creativecommons.org/licenses/by/4.0/, as recorded in the arXiv licence metadata for this exact version and shown on the abstract page https://arxiv.org/abs/2606.06403v1. Full licence notice: \texttt{licenses/arxiv-2606.06403-CC-BY-4.0.txt}.\par\smallskip

\noindent\textit{Editorial scope.} Counted unit: the article's proposition proposition (source0.tex lines 809--836). The article's complete original proof of it is delivered with it (source0.tex lines 838--957, one \texttt{proof} environment of 120 lines). No mathematics is written, repaired or re-derived: the statement and the proof are the article's own lines, in the article's own notation, and no retained line is substituted or re-flowed.  No further source-local context is needed: every cross-reference of every retained line resolves inside the two delivered spans.  Nothing was omitted from any retained span, so the package carries no omission record.  No retained line contains a citation, so the wrapper carries no bibliography and no bibliography entry.  No retained line contains a figure environment, an image-inclusion command or drawing code, so no image is delivered and the package declares no asset dependency.  Editorial formatting only, in addition: an AMS \texttt{amsart} preamble that restates the article's own declarations and macros used by the retained lines, namely the environments \texttt{proposition} on separate counters, together with the standard packages those lines need.  The retained environments print in this document as Proposition 1 in order of appearance, whereas the article prints the same items with the numbering of its own full document.  The complete native source of the article as distributed in the arXiv e-print for this exact version is bundled unchanged as \texttt{sources/202187/source0.tex}.
\begin{proposition}[Spin-Fourier residues in the square family]
Let $\ell\geq 5$ be odd.  For
\begin{equation}
p=\ell^2,
\qquad
q_1=\ell-1,
\qquad
q_2=2\ell-1,
\qquad
k=\ell-1,
\end{equation}
one has
\begin{equation}
m_k(p,q_1)=m_k(p,q_2)=(2,2).
\end{equation}
However,
\begin{equation}
\Delta_k(p,q_1;u)=A_{\ell-2}(u)-A_{\ell}(u),
\end{equation}
whereas
\begin{equation}
\Delta_k(p,q_2;u)=A_{2\ell-1}(u)-A_{\ell}(u).
\end{equation}
In particular,
\begin{equation}
\Delta_k(p,q_1;u)\neq \Delta_k(p,q_2;u).
\end{equation}
\end{proposition}

\begin{proof}
At level $k=\ell-1$, the contributing affine spin-lattice points satisfy
\begin{equation}
|a|+|b|\leq 2\ell
\end{equation}
and
\begin{equation}
a+qb\equiv0\pmod{\ell^2}.
\end{equation}
First take $q=q_1=\ell-1$.  The congruence becomes
\begin{equation}
a+(\ell-1)b\equiv0\pmod{\ell^2}.
\end{equation}
Modulo $\ell$, this implies
\begin{equation}
a-b\equiv0\pmod{\ell}.
\end{equation}
Since $a$ and $b$ are odd, the difference $a-b$ is even.  Under the bound $|a|+|b|\leq2\ell$, the only possible cases are
\begin{equation}
a=b,
\qquad
a=b+2\ell,
\qquad
a=b-2\ell.
\end{equation}
If $a=b$, then
\begin{equation}
a+(\ell-1)b=\ell a.
\end{equation}
Divisibility by $\ell^2$ forces $a$ to be a multiple of $\ell$.  Since $|a|\leq\ell$ and $a$ is odd, this gives
\begin{equation}
(a,b)=(\ell,\ell)
\qquad\text{or}\qquad
(a,b)=(-\ell,-\ell).
\end{equation}
If $a=b+2\ell$, then
\begin{equation}
a+(\ell-1)b=\ell(b+2).
\end{equation}
Divisibility by $\ell^2$ gives
\begin{equation}
b+2\equiv0\pmod{\ell}.
\end{equation}
The bound gives
\begin{equation}
b=-(\ell+2)
\qquad\text{and}\qquad
a=\ell-2.
\end{equation}
Thus one obtains
\begin{equation}
(a,b)=(\ell-2,-\ell-2).
\end{equation}
The case $a=b-2\ell$ similarly gives
\begin{equation}
(a,b)=(-\ell+2,\ell+2).
\end{equation}
Therefore the contributing first Spin-Fourier weights for $q_1$ are
\begin{equation}
\pm\ell
\qquad\text{and}\qquad
\pm(\ell-2).
\end{equation}
The points $(\ell,\ell)$ and $(-\ell,-\ell)$ have even sign-counting parity and contribute to the negative branch at this level.  The points $(\ell-2,-\ell-2)$ and $(-\ell+2,\ell+2)$ have odd sign-counting parity and contribute to the positive branch.  Hence
\begin{equation}
\Delta_k(p,q_1;u)=A_{\ell-2}(u)-A_{\ell}(u).
\end{equation}

Now take $q=q_2=2\ell-1$.  The congruence is
\begin{equation}
a+(2\ell-1)b\equiv0\pmod{\ell^2}.
\end{equation}
Again, modulo $\ell$, this implies
\begin{equation}
a-b\equiv0\pmod{\ell}.
\end{equation}
The same three cases occur.  If $a=b$, then
\begin{equation}
a+(2\ell-1)b=2\ell a.
\end{equation}
Since $\ell$ is odd, divisibility by $\ell^2$ forces $a$ to be a multiple of $\ell$.  Hence
\begin{equation}
(a,b)=(\ell,\ell)
\qquad\text{or}\qquad
(a,b)=(-\ell,-\ell).
\end{equation}
If $a=b+2\ell$, then
\begin{equation}
a+(2\ell-1)b=2\ell(b+1).
\end{equation}
Divisibility by $\ell^2$ gives
\begin{equation}
b+1\equiv0\pmod{\ell}.
\end{equation}
The bound gives
\begin{equation}
b=-1
\qquad\text{and}\qquad
a=2\ell-1.
\end{equation}
Thus one obtains
\begin{equation}
(a,b)=(2\ell-1,-1).
\end{equation}
The case $a=b-2\ell$ similarly gives
\begin{equation}
(a,b)=(-(2\ell-1),1).
\end{equation}
Therefore the contributing first Spin-Fourier weights for $q_2$ are
\begin{equation}
\pm\ell
\qquad\text{and}\qquad
\pm(2\ell-1).
\end{equation}
The $\pm\ell$ pair contributes to the negative branch, while the $\pm(2\ell-1)$ pair contributes to the positive branch.  Hence
\begin{equation}
\Delta_k(p,q_2;u)=A_{2\ell-1}(u)-A_{\ell}(u).
\end{equation}
Both examples have one positive reduced weight pair and one negative reduced weight pair, so their scalar data are both $(2,2)$.  Since $A_{\ell-2}(u)$ and $A_{2\ell-1}(u)$ are distinct reduced Spin-Fourier generators, the two residues are different.
\end{proof}

\end{document}
